\section{Conception de la watchapp (montre)}
\label{sec:montre}

\subsection{Composants}

\begin{table}[H]
  \centering\small
  \begin{tabularx}{\linewidth}{@{}l>{\raggedright\arraybackslash}p{30mm}X>{\raggedright\arraybackslash}p{30mm}@{}}
    \toprule
    \textbf{Réf.} & \textbf{Composant} & \textbf{Responsabilité} & \textbf{Exigences} \\
    \midrule
    M-01 & Application & Cycle de vie, pile de fenêtres, ouverture de la messagerie, affichage immédiat de la dernière situation. & \refexig{TPB-NAV-001}, \refexig{TPB-NAV-005}, \refexig{TPB-ACT-001} \\
    M-02 & Messagerie & Ouvre AppMessage, envoie les demandes l'une après l'autre (attente de l'accusé d'envoi, un réessai), reçoit les vues, ignore les messages d'une demande périmée. & \refexig{TPB-ACT-009} \\
    M-03 & Magasin & Garde en mémoire la vue de l'écran courant~; écrit dans le stockage la dernière liste des pièces avec son heure. & \refexig{TPB-INF-001}, \refexig{TPB-INF-002} \\
    M-04 & Plateforme & Table de constantes par plateforme (§\ref{sec:plateformes}). Seul endroit où l'on teste la plateforme. & \refexig{TPB-PLT-001} à \refexig{TPB-PLT-007} \\
    M-05 & Thème & Couleurs, motifs noir et blanc, tranches de température (§\ref{sec:theme}). & \refexig{TPB-AFF-001} à \refexig{TPB-AFF-004} \\
    M-06 & Moteur de liste & Une seule liste pour tous les écrans de liste~: rectangulaire, ou ronde sur trois lignes avec la sélection au centre. Défilement, sélection, sous-titres, pastilles. & \refexig{TPB-NAV-002}, \refexig{TPB-PLT-004} \\
    M-07 & Fenêtres & Une fenêtre par écran (§\ref{sec:fenetres}). Elles ne connaissent ni Tado ni le protocole~: elles lisent la vue du magasin et émettent des demandes. & \refexig{TPB-PIE-001} à \refexig{TPB-ALE-005} \\
    M-08 & Textes & Ressources en français et en anglais, mise en forme des températures et des heures. & \refexig{TPB-AFF-005}, \refexig{TPB-AFF-006} \\
    M-09 & Barre d'actions & Barre latérale des écrans rectangulaires (30~px sur Time~2, 24~px ailleurs). Absente des écrans ronds. & \refexig{TPB-PLT-005} \\
    M-10 & Bandeau d'état & Âge des données, absence de connexion, attente de réponse. & \refexig{TPB-ACT-002}, \refexig{TPB-UTI-005} \\
    \bottomrule
  \end{tabularx}
  \caption{Composants de la watchapp}
\end{table}

\subsection{Adaptation aux plateformes}
\label{sec:plateformes}

\begin{table}[H]
  \centering\small
  \begin{tabularx}{\linewidth}{@{}l>{\raggedright\arraybackslash}p{34mm}>{\raggedright\arraybackslash}p{22mm}>{\raggedright\arraybackslash}p{18mm}X@{}}
    \toprule
    \textbf{Plateforme} & \textbf{Montres} & \textbf{Écran} & \textbf{Stockage} & \textbf{Particularités de présentation} \\
    \midrule
    emery & Pebble Time~2 & 200\,×\,228, couleur & 1~Mo & Référence. Barre d'actions de 30~px, sous-titres visibles, légende complète de l'historique, douze colonnes. \\
    basalt & Pebble Time, Time Steel & 144\,×\,168, couleur & 4~Ko & Barre d'actions de 24~px, sous-titres abrégés, noms de pièces abrégés dans l'historique. \\
    chalk & Pebble Time Round & 180\,×\,180, rond, couleur & 4~Ko & Trois lignes, sélection au centre, pas de barre d'actions, historique sur huit colonnes. \\
    gabbro & Pebble Round~2 & 260\,×\,260, rond, couleur & 1~Mo & Comme chalk, avec les textes et douze colonnes d'emery. \\
    diorite & Pebble~2 & 144\,×\,168, N\&B & 4~Ko & Vidéo inversée, contours, motifs. \\
    flint & Pebble~2 Duo & 144\,×\,168, N\&B & 1~Mo & Comme diorite (rendu identique). \\
    \bottomrule
  \end{tabularx}
  \caption{Plateformes et choix de présentation}
\end{table}

La table de constantes (M-04) porte, par plateforme~: hauteurs de ligne et de sous-titre,
largeur de barre d'actions, nombre de colonnes d'historique, présence de la légende, indicateur
\enquote{rond}, indicateur \enquote{couleur}. Elle se remplit en relevant les valeurs sur les
maquettes (§\ref{sec:maquettes}). On teste \texttt{PBL\_ROUND}, \texttt{PBL\_COLOR} et
\texttt{PBL\_DISPLAY\_WIDTH} plutôt que le nom de la plateforme, qui n'est utilisé que
pour la table.

\subsection{Thème}
\label{sec:theme}

Le thème \enquote{Tado tuiles} choisi sur les maquettes donne un fond noir, une barre de titre
orange, une sélection ambre, et des tuiles de couleur selon la température (\refexig{TPB-AFF-001}).

\begin{table}[H]
  \centering\small
  \begin{tabularx}{\linewidth}{@{}>{\raggedright\arraybackslash}p{34mm}>{\raggedright\arraybackslash}p{30mm}>{\raggedright\arraybackslash}p{50mm}X@{}}
    \toprule
    \textbf{Usage} & \textbf{Teinte} & \textbf{Constante du SDK (à vérifier)} & \textbf{Noir et blanc} \\
    \midrule
    Pièce sous 18~°C        & \texttt{\#00AAFF} & \texttt{GColorVividCerulean} & ligne blanche \\
    18 à 19,5~°C            & \texttt{\#55FFFF} & \texttt{GColorCeleste}       & ligne blanche \\
    19,5 à 21~°C            & \texttt{\#FFAA00} & \texttt{GColorChromeYellow}  & ligne blanche \\
    21~°C et plus           & \texttt{\#FF5500} & \texttt{GColorOrange}        & ligne blanche \\
    Pièce qui chauffe       & un à trois traits de chaleur et pourcentage & \texttt{GColorOrange}    & ligne inversée \\
    Titre                   & \texttt{\#FF5500} & \texttt{GColorOrange}        & bandeau noir \\
    Historique, niveaux 1 à 3 & \texttt{\#AA5500}, \texttt{\#FF5500}, \texttt{\#FFAA00} & \texttt{GColorWindsorTan}, \texttt{GColorOrange}, \texttt{GColorChromeYellow} & points, damier, plein \\
    \bottomrule
  \end{tabularx}
  \caption{Couleurs, toutes dans la palette de 64 couleurs (\refexig{TPB-AFF-002})}
\end{table}

Le niveau de chauffe se lit sur les traits de chaleur de la liste, comme dans l'application Tado~: un trait
jusqu'à 33~\%, deux jusqu'à 66~\%, trois au-delà (seuils à confirmer sur l'application Android)~; le
pourcentage est écrit à côté (\refexig{TPB-AFF-003}).

Les seuils de température sont comparés dans l'unité du compte~: en degrés Fahrenheit, ils sont
convertis une fois (64,4~; 67,1~; 69,8~°F) par le pkjs, qui envoie la \emph{tranche} avec la
température. La watchapp ne compare donc jamais de seuils (DEC-02).

\subsection{Fenêtres et navigation}
\label{sec:fenetres}

\begin{table}[H]
  \centering\small
  \begin{tabularx}{\linewidth}{@{}>{\raggedright\arraybackslash}p{27mm}>{\raggedright\arraybackslash}p{22mm}X@{}}
    \toprule
    \textbf{Fenêtre} & \textbf{Parent} & \textbf{Boutons} \\
    \midrule
    1 Pièces & aucun & Haut/Bas~: sélection. Sélection~: détail. Sélection long~: menu. Retour~: quitter. \\
    2 Détail & Pièces & Haut/Bas (barre d'actions)~: pièce précédente/suivante. Sélection~: pilotage. Sélection long~: bien-être. Retour~: pièces. \\
    3 Pilotage & Détail & Haut/Bas~: sélection. Sélection~: action (forcer, rejouer, programmation, éteindre). Retour~: détail. \\
    4 Consigne & Pilotage & Haut/Bas~: $\pm 0{,}5$~° ($\pm 1$~°F). Sélection~: valider, ouvre la durée. Retour~: abandonner sans effet. \\
    5 Durée & Consigne & Haut/Bas~: sélection. Sélection~: envoyer la commande. Retour~: consigne. \\
    6 Menu & Pièces & Haut/Bas~: sélection. Sélection~: ouvrir. Retour~: pièces. \\
    7 Historique & Menu & Retour seulement. Haut/Bas~: défilement s'il y a plus de pièces que de lignes. \\
    8 Présence & Menu & Haut/Bas~: sélection. Sélection~: imposer le mode. Retour~: menu. \\
    9 Raccourcis & Menu & Haut/Bas~: sélection. Sélection~: confirmation, puis commande. Retour~: menu. \\
    10 Bien-être & Détail & Haut/Bas~: pièce précédente/suivante. Retour~: détail. \\
    11 Alertes & Menu & Haut/Bas~: sélection. Retour~: menu. \\
    12 Liaison & Pièces & Retour~: annuler la liaison. \\
    \bottomrule
  \end{tabularx}
  \caption{Fenêtres et fonction des boutons (\refexig{TPB-NAV-002} à \refexig{TPB-NAV-005})}
  \label{tab:fenetres}
\end{table}

Une fenêtre n'existe que pendant qu'elle est empilée~: elle est créée à l'ouverture, détruite
à la fermeture, et ne garde rien que la vue du magasin ne sache redonner. Les données de chaque écran sont
rechargées à son affichage (ouverture, retour d'une fenêtre qui le recouvrait)~: le téléphone relit alors les
pièces chez Tado et les renvoie avec la réponse, sous une même séquence.

\paragraph{Plateforme abandonnée.} aplite n'est pas prise en charge (DEC-14)~: son code occupait 21~Ko de ses
24~Ko de mémoire de travail.

\subsection{Données conservées sur la montre}

Le magasin écrit, après chaque liste de pièces reçue, un enregistrement unique~:
version du format (1 octet), heure d'obtention (4), nombre de pièces (1), puis pour chaque pièce son
nom (20), son abréviation (4), sa température mesurée et sa consigne (2 chacune, en dixièmes), sa
puissance (1), ses indicateurs (1) et sa tranche de couleur (1), soit 34 octets par pièce. Seize
pièces tiennent dans 600 octets, loin des 4~Ko de la plus petite plateforme (C7). Rien d'autre
n'est écrit~: pas de jeton, pas d'identifiant Tado.
